% GENERATED FILE. Edit canonical lessons, then run npm run generate:books.
% canonical-source-hash: fnv1a64:8743f3ff7d763789
% canonical-lessons: TE-C316-vellanu, TE-C316-cestunnappudu, TE-C316-cadivevadini, TE-R316-the-past-three-ways

\chapter{Was Doing, Did, Used To}
\label{ch:te-was-doing-did-used-to}
\hlchaptermodality{316}

\begin{chapteropening}
\textbf{By the end of this chapter:} \emph{I can say what happened in the plain past (veḷḷānu), set the scene with cēstunnappuḍu, and say what used to happen with cadivēvāḍini and uṇḍēdi.}

Complete TE-R316-the-past-three-ways: tell what happened, what was going on and what used to happen, in the plain past, -tunnappuḍu and -ēvāḍini.
\end{chapteropening}

\section[\texorpdfstring{\te{వెళ్ళాను} (veḷḷānu)}{వెళ్ళాను (veḷḷānu)}]{\te{వెళ్ళాను} (veḷḷānu) — the plain past: what happened}

\label{lesson:TE-C316-vellanu}



\begin{quote}
Before the new one: say \textquotedblleft{}I go\textquotedblright{}, or \textquotedblleft{}I will go\textquotedblright{}. (\emph{nēnu veḷtānu}.) Then the last two people you met: the watchman and the milkman. (\emph{kāpalādāru}, \emph{pālavāḍu}.)

\textbf{\te{వెళ్తాను}} says what happens, or what will. A story has already happened. It needs the past.
\end{quote}

\subsection*{You'll want to know: the plain past}
\begin{quote}
\textbf{\te{నిన్న} \te{నేను} \te{పట్టణానికి} \te{వెళ్ళాను}.}
\end{quote}

\begin{quote}
\emph{ninna nēnu paṭṭaṇāniki veḷḷānu.} — \textquotedblleft{}Yesterday I went to the town.\textquotedblright{}
\end{quote}

\textbf{\te{వెళ్ళాను}} is \textquotedblleft{}I went\textquotedblright{}. Where the form you know has \textbf{-\te{తా}-} (\textbf{\te{వెళ్తాను}}), the past has \textbf{-\te{ఆ}-}: \textbf{\te{వెళ్ళ}} + \textbf{\te{ఆను}}. The person ending at the end works the way it does in \textbf{\te{వెళ్తాను}}:

\begin{itemize}
\raggedright
  \item \textbf{\te{నేను} \te{వెళ్ళాను}} (\emph{nēnu veḷḷānu}) — I went
  \item \textbf{\te{నువ్వు} \te{వెళ్ళావు}} (\emph{nuvvu veḷḷāvu}) — you went
  \item \textbf{\te{అతను} \te{వెళ్ళాడు}} (\emph{atanu veḷḷāḍu}) — he went
  \item \textbf{\te{ఆమె} \te{వెళ్ళింది}} (\emph{āme veḷḷindi}) — she went, and \textquotedblleft{}it went\textquotedblright{}
  \item \textbf{\te{మేము} \te{వెళ్ళాము}} (\emph{mēmu veḷḷāmu}) — we went
  \item \textbf{\te{వాళ్ళు} \te{వెళ్ళారు}} (\emph{vāḷḷu veḷḷāru}) — they went; with \textbf{\te{మీరు}}, \textquotedblleft{}you went\textquotedblright{}, said with respect
\end{itemize}

Only \textquotedblleft{}she\textquotedblright{} and \textquotedblleft{}it\textquotedblright{} step out of line: \textbf{\te{వెళ్ళింది}}, with \textbf{-\te{ఇంది}}.

\subsection*{You'll want to know: the stems that change}
Many verbs take the past on their own stem. A \textbf{-\te{ఉ}-} before the last consonant turns to \textbf{-\te{ఇ}-}: \textbf{\te{చదువు}} $\to$ \textbf{\te{చదివాను}} (\emph{cadivānu}), \textquotedblleft{}I read\textquotedblright{}; \textbf{\te{నడుచు}} $\to$ \textbf{\te{నడిచాను}} (\emph{naḍicānu}), \textquotedblleft{}I walked\textquotedblright{}.

A few of the commonest verbs change more, and they are worth learning as pairs:

\begin{itemize}
\raggedright
  \item \textbf{\te{వచ్చు}} $\to$ \textbf{\te{వచ్చాను}} (\emph{vaccānu}), \textquotedblleft{}I came\textquotedblright{}
  \item \textbf{\te{చూడు}} $\to$ \textbf{\te{చూశాను}} (\emph{cūśānu}), \textquotedblleft{}I saw\textquotedblright{}; she saw: \textbf{\te{చూసింది}} (\emph{cūsindi})
  \item \textbf{\te{చేయు}} $\to$ \textbf{\te{చేశాను}} (\emph{cēśānu}), \textquotedblleft{}I did\textquotedblright{}; she did: \textbf{\te{చేసింది}} (\emph{cēsindi})
  \item \textbf{\te{తిను}} $\to$ \textbf{\te{తిన్నాను}} (\emph{tinnānu}), \textquotedblleft{}I ate\textquotedblright{}
  \item \textbf{\te{ఇవ్వు}} $\to$ \textbf{\te{ఇచ్చాను}} (\emph{iccānu}), \textquotedblleft{}I gave\textquotedblright{}
\end{itemize}

\subsection*{Your turn}
Say these aloud:

\begin{itemize}
\raggedright
  \item \textquotedblleft{}Yesterday I went to the town\textquotedblright{} — \emph{ninna nēnu paṭṭaṇāniki veḷḷānu}
  \item \textquotedblleft{}She came in the evening\textquotedblright{} — \emph{āme sāyantraṁ vaccindi}
  \item \textquotedblleft{}We drank tea\textquotedblright{} — \emph{mēmu ṭī tāgāmu}
  \item \textquotedblleft{}I read a book. After that I had a meal.\textquotedblright{} — \emph{nēnu oka pustakaṁ cadivānu. ā taruvāta bhōjanaṁ cēśānu.}
\end{itemize}

\begin{tcolorbox}[breakable,colback=teal!4,colframe=teal!35!black,title={Before you move on}]
\textquotedblleft{}I went\textquotedblright{}? (\textbf{\te{వెళ్ళాను}}.) \textquotedblleft{}He came\textquotedblright{}? (\textbf{\te{వచ్చాడు}}.) \textquotedblleft{}She saw\textquotedblright{}? (\textbf{\te{చూసింది}}.) \textquotedblleft{}I ate\textquotedblright{}? (\textbf{\te{తిన్నాను}}.) Which person takes \textbf{-\te{ఇంది}} instead of \textbf{-\te{ఆ}-}? (\textbf{She, and it.})
\end{tcolorbox}
\section[\texorpdfstring{\te{చేస్తున్నప్పుడు} (cēstunnappuḍu)}{చేస్తున్నప్పుడు (cēstunnappuḍu)}]{\te{చేస్తున్నప్పుడు} (cēstunnappuḍu) — the scene and the event}

\label{lesson:TE-C316-cestunnappudu}



\begin{quote}
Say \textquotedblleft{}I am eating now\textquotedblright{}. (\emph{nēnu ippuḍu tiṇṭunnānu}.) Then \textquotedblleft{}I ate\textquotedblright{}. (\emph{tinnānu}.)

The first is going on. The second is done. A told story needs both, one inside the other.
\end{quote}

\subsection*{You'll want to know: -\te{తున్నప్పుడు}, while it was going on}
\begin{quote}
\textbf{\te{నేను} \te{భోజనం} \te{చేస్తున్నప్పుడు} \te{ఫోను} \te{వచ్చింది}.}
\end{quote}

\begin{quote}
\emph{nēnu bhōjanaṁ cēstunnappuḍu phōnu vaccindi.} — \textquotedblleft{}While I was having a meal, a call came.\textquotedblright{}
\end{quote}

Take \textbf{\te{చేస్తున్నాను}}, \textquotedblleft{}I am doing\textquotedblright{}, drop the person ending, and add \textbf{\te{అప్పుడు}}, \textquotedblleft{}then, at that time\textquotedblright{}: \textbf{\te{చేస్తున్న}} + \textbf{\te{అప్పుడు}} $\to$ \textbf{\te{చేస్తున్నప్పుడు}}, \textquotedblleft{}at the time (I) was doing\textquotedblright{}.

That clause is the \textbf{scene}. The thing that cut into it is the \textbf{event}, in the plain past: \textbf{\te{వచ్చింది}}. Telugu says a call \textquotedblleft{}came\textquotedblright{}, where English says the phone rang.

\begin{quote}
\textbf{\te{మేము} \te{నడుస్తున్నప్పుడు} \te{అకస్మాత్తుగా} \te{వర్షం} \te{పడింది}.}
\end{quote}

\begin{quote}
\emph{mēmu naḍustunnappuḍu akasmāttugā varṣaṁ paḍindi.} — \textquotedblleft{}We were walking when suddenly it rained.\textquotedblright{}
\end{quote}

\textbf{\te{వర్షం}} (\emph{varṣaṁ}, \textquotedblleft{}rain\textquotedblright{}) with \textbf{\te{పడింది}} (\emph{paḍindi}, \textquotedblleft{}fell\textquotedblright{}): rain \textquotedblleft{}falls\textquotedblright{} in Telugu.

\begin{grammarlens}[title={Grammar Lens: 'was doing' on its own}]
Telugu needs no separate tense for \textquotedblleft{}was doing\textquotedblright{}. The same form that says \textquotedblleft{}I am reading\textquotedblright{} says \textquotedblleft{}I was reading\textquotedblright{} once a past time sets the scene:

\begin{quote}
\textbf{\te{నిన్న} \te{రాత్రి} \te{తొమ్మిది} \te{గంటలకు} \te{నేను} \te{పుస్తకం} \te{చదువుతున్నాను}.}
\end{quote}

\begin{quote}
\emph{ninna rātri tommidi gaṇṭalaku nēnu pustakaṁ caduvutunnānu.} — \textquotedblleft{}At nine last night I was reading a book.\textquotedblright{}
\end{quote}

The past is carried by \textbf{\te{నిన్న} \te{రాత్రి}}, not by the verb. \textbf{\te{ఉన్నాను}}, the \textquotedblleft{}am\textquotedblright{} inside \textbf{\te{చదువుతున్నాను}}, works the same way:

\begin{quote}
\textbf{\te{నిన్న} \te{నేను} \te{ఇంట్లో} \te{ఉన్నాను}.}
\end{quote}

\begin{quote}
\emph{ninna nēnu iṇṭlō unnānu.} — \textquotedblleft{}Yesterday I was at home.\textquotedblright{}
\end{quote}

\textbf{\te{ఇంట్లో}} (\emph{iṇṭlō}, \textquotedblleft{}at home\textquotedblright{}) is the \textbf{\te{ఇంటి}-} of \textbf{\te{ఇంటికి}} with \textbf{-\te{లో}}, \textquotedblleft{}in\textquotedblright{}.
\end{grammarlens}

\subsection*{Your turn}
Say these aloud:

\begin{itemize}
\raggedright
  \item \textquotedblleft{}While I was having a meal, a call came\textquotedblright{} — \emph{nēnu bhōjanaṁ cēstunnappuḍu phōnu vaccindi}
  \item \textquotedblleft{}While the watchman was sleeping, the milkman came\textquotedblright{} — \emph{kāpalādāru nidrapōtunnappuḍu pālavāḍu vaccāḍu}
  \item \textquotedblleft{}At nine last night I was reading a book\textquotedblright{} — \emph{ninna rātri tommidi gaṇṭalaku nēnu pustakaṁ caduvutunnānu}
\end{itemize}

\begin{tcolorbox}[breakable,colback=teal!4,colframe=teal!35!black,title={Before you move on}]
Which form paints the scene? (\textbf{-\te{తున్నప్పుడు}}, or \textbf{-\te{తున్నాను}} with a past time.) Which form drops in the event? (The plain past: \textbf{\te{వచ్చింది}}, \textbf{\te{పడింది}}.) \textquotedblleft{}While I was eating\textquotedblright{}? (\textbf{\te{నేను} \te{తింటున్నప్పుడు}}, \emph{nēnu tiṇṭunnappuḍu}.)
\end{tcolorbox}
\section[\texorpdfstring{\te{చదివేవాడిని} (cadivēvāḍini)}{చదివేవాడిని (cadivēvāḍini)}]{\te{చదివేవాడిని} (cadivēvāḍini) — what used to happen}

\label{lesson:TE-C316-cadivevadini}



\begin{quote}
Say \textquotedblleft{}While I was having a meal, a call came.\textquotedblright{} (\emph{nēnu bhōjanaṁ cēstunnappuḍu phōnu vaccindi}.)

That happened once. Some things happened many times: what you did every day as a child.
\end{quote}

\subsection*{You'll want to know: -\te{ఏవాడిని}, 'used to'}
\begin{quote}
\textbf{\te{బాల్యంలో} \te{నేను} \te{ప్రతి} \te{రోజు} \te{ఒక} \te{కథ} \te{చదివేవాడిని}.}
\end{quote}

\begin{quote}
\emph{bālyaṁlō nēnu prati rōju oka katha cadivēvāḍini.} — \textquotedblleft{}In my childhood I used to read a story every day.\textquotedblright{}
\end{quote}

\textbf{\te{చదివేవాడిని}} is \textquotedblleft{}I used to read\textquotedblright{}, said by a man. A woman says \textbf{\te{చదివేదాన్ని}} (\emph{cadivēdānni}). Taken apart, it is the \textbf{-\te{ఏ}} form of the verb (\textbf{\te{చదివే}}, \textquotedblleft{}who reads\textquotedblright{}) with \textbf{\te{వాడు}}, \textquotedblleft{}that man\textquotedblright{}, or \textbf{\te{దాని}}, \textquotedblleft{}that one\textquotedblright{}, and the \textbf{-\te{ని}} of \textquotedblleft{}I\textquotedblright{}: roughly \textquotedblleft{}I was one who reads\textquotedblright{}.

\begin{itemize}
\raggedright
  \item \textbf{\te{నేను} \te{చదివేవాడిని}} (\emph{nēnu cadivēvāḍini}) / \textbf{\te{చదివేదాన్ని}} (\emph{cadivēdānni}) — I used to read
  \item \textbf{\te{అతను} \te{చదివేవాడు}} (\emph{atanu cadivēvāḍu}) — he used to read
  \item \textbf{\te{ఆమె} \te{చదివేది}} (\emph{āme cadivēdi}) — she used to read
  \item \textbf{\te{మేము} \te{చదివేవాళ్ళం}} (\emph{mēmu cadivēvāḷḷaṁ}) — we used to read
  \item \textbf{\te{వాళ్ళు} \te{చదివేవాళ్ళు}} (\emph{vāḷḷu cadivēvāḷḷu}) — they used to read
\end{itemize}

The same \textbf{-\te{ఏ}} sits on any verb: \textbf{\te{వెళ్ళేవాడిని}}, \textquotedblleft{}I used to go\textquotedblright{}; \textbf{\te{తినేవాడిని}}, \textquotedblleft{}I used to eat\textquotedblright{}.

\subsection*{You'll want to know: \te{ఉండేది}, 'there used to be'}
\begin{quote}
\textbf{\te{మా} \te{గ్రామంలో} \te{ఒక} \te{చిన్న} \te{అంగడి} \te{ఉండేది}.}
\end{quote}

\begin{quote}
\emph{mā grāmaṁlō oka cinna aṅgaḍi uṇḍēdi.} — \textquotedblleft{}There used to be a small shop in our village.\textquotedblright{}
\end{quote}

\textbf{\te{ఉండేది}} is the same form of \textbf{\te{ఉండు}}, \textquotedblleft{}to be\textquotedblright{}: \textquotedblleft{}it used to be there\textquotedblright{}. Set it against the present:

\begin{quote}
\textbf{\te{ఇప్పుడు} \te{అక్కడ} \te{ఒక} \te{పెద్ద} \te{హోటలు} \te{ఉంది}.}
\end{quote}

\begin{quote}
\emph{ippuḍu akkaḍa oka pedda hōṭalu undi.} — \textquotedblleft{}Now there is a big restaurant there.\textquotedblright{}
\end{quote}

A story about the old days leans on \textbf{\te{ఉండేది}} for the background.

\subsection*{Your turn}
Say these aloud:

\begin{itemize}
\raggedright
  \item \textquotedblleft{}In my childhood I used to read a story every day\textquotedblright{} — \emph{bālyaṁlō nēnu prati rōju oka katha cadivēvāḍini}
  \item \textquotedblleft{}A woman says: I used to go to school every day\textquotedblright{} — \emph{nēnu prati rōju baḍiki veḷḷēdānni}
  \item \textquotedblleft{}My grandmother used to cook every evening\textquotedblright{} — \emph{mā ammamma prati sāyantraṁ vaṇḍēdi}
\end{itemize}

\begin{tcolorbox}[breakable,colback=teal!4,colframe=teal!35!black,title={Before you move on}]
\textquotedblleft{}I used to go\textquotedblright{}, said by a man? (\textbf{\te{వెళ్ళేవాడిని}}.) By a woman? (\textbf{\te{వెళ్ళేదాన్ని}}.) \textquotedblleft{}There used to be\textquotedblright{}? (\textbf{\te{ఉండేది}}.) Which vowel marks the habit? (\textbf{-\te{ఏ}}, as in \textbf{\te{చదివే}}.)
\end{tcolorbox}
\section[(review)]{Review — the past three ways}

\label{lesson:TE-R316-the-past-three-ways}



\begin{quote}
Say \textquotedblleft{}Yesterday I went to the town.\textquotedblright{} (\emph{ninna nēnu paṭṭaṇāniki veḷḷānu}.) Then \textquotedblleft{}she saw\textquotedblright{} and \textquotedblleft{}I ate\textquotedblright{}. (\emph{cūsindi}, \emph{tinnānu}.)
\end{quote}

\subsection*{Your turn}
Say these aloud:

\begin{itemize}
\raggedright
  \item We drank tea, and after that we saw a film.
  \item While I was reading a book, Mother called.
  \item At nine last night I was at home.
  \item In my childhood I used to walk to school every day.
  \item There used to be a small shop in our village. Now there is a big restaurant.
  \item one thing you were doing yesterday evening, with -\te{తున్నప్పుడు}, and one thing that cut into it
\end{itemize}

\begin{tcolorbox}[breakable,colback=teal!4,colframe=teal!35!black,title={Before you move on}]
Which past is the event? (The plain past: \textbf{\te{వెళ్ళాను}}, \textbf{\te{వచ్చింది}}.) Which is the scene? (\textbf{-\te{తున్నప్పుడు}}, or \textbf{-\te{తున్నాను}} with a past time.) Which is the habit? (\textbf{-\te{ఏవాడిని}}, \textbf{-\te{ఏదాన్ని}}.) \textquotedblleft{}There used to be\textquotedblright{}? (\textbf{\te{ఉండేది}}.)
\end{tcolorbox}
